% =====================================================================
% Vlastní makra
%
% Soubor předpokládá balíčky z packages.tex a barvy ze style.tex.
% =====================================================================

% --- Zvýrazňování ----------------------------------------------------
\newcommand{\markred}[1]{\textcolor{lightred}{#1}}
\newcommand{\markdarkred}[1]{\textcolor{darkred}{#1}}
\newcommand{\markgreen}[1]{\textcolor{darkgreen}{#1}}
\newcommand{\markblue}[1]{\textcolor{lightblue}{#1}}
\newcommand{\markdarkblue}[1]{\textcolor{darkblue}{#1}}

\newcommand{\cmark}{\markgreen{\ding{51}}}
\newcommand{\xmark}{\markred{\ding{55}}}

% Poznámka pro autora. Ve výsledném textu je červená, aby nešla přehlédnout.
\newcommand{\todo}[1]{\textcolor{darkred}{\textbf{(TODO: #1)}}}

% --- Číselné množiny a základní značení ------------------------------
\newcommand{\R}{\mathbb{R}}
\newcommand{\C}{\mathbb{C}}
\newcommand{\N}{\mathbb{N}}
\newcommand{\Z}{\mathbb{Z}}
\newcommand{\Q}{\mathbb{Q}}

\renewcommand{\emptyset}{\varnothing}

\newcommand{\set}[1]{\left\{#1\right\}}
\newcommand{\powset}{\mathcal{P}}
\newcommand{\code}[1]{\langle#1\rangle}
\newcommand{\T}[1]{#1^\top}                    % transpozice
\newcommand{\compl}[1]{#1^\complement}         % doplněk
\newcommand{\mapping}[3]{#1\colon #2 \rightarrow #3}
\newcommand{\mapto}[3]{#1\colon #2 \mapsto #3}
\newcommand{\bigO}{\mathcal{O}}
% Jméno \binary si zabírá newtxmath i newpxmath, proto \binrep.
\newcommand{\binrep}[1]{\ensuremath{\left\langle #1\right\rangle}_B}

\DeclareMathOperator*{\argmin}{arg\,min}
\DeclareMathOperator*{\argmax}{arg\,max}
\DeclareMathOperator{\mspan}{span}
\DeclareMathOperator{\mrank}{rank}
\DeclareMathOperator{\id}{id}
\DeclareMathOperator{\dom}{dom}
\DeclareMathOperator{\rng}{rng}

% --- Statistika ------------------------------------------------------
\newcommand{\average}[1]{\overline{#1}}
\DeclareMathOperator{\median}{med}
\DeclareMathOperator{\Var}{var}
\DeclareMathOperator{\Cov}{cov}
\DeclareMathOperator{\MSE}{MSE}
\DeclareMathOperator{\SSE}{SSE}

% --- Obrázky ---------------------------------------------------------
% Obrázek na střed s popiskem a návěštím. Výška se omezuje shora, aby
% se ani vysoký obrázek nevytlačil na vlastní stranu.
% #1 = maximální výška (volitelně), #2 = soubor, #3 = popisek, #4 = label
\newcommand{\obrazek}[4][0.45\textheight]{%
    \begin{figure}[H]
        \centering
        \includegraphics[max height=#1, max width=\textwidth]{#2}
        \caption{#3}
        \label{#4}
    \end{figure}%
}

% Ikona sázená do řádku textu, zarovnaná na účaří.
\newcommand{\icon}[1]{%
    \raisebox{-0.15\height}{\includegraphics[height=1em]{#1}}%
}

% --- Rámečky ---------------------------------------------------------
% Čtyři úrovně důrazu se stejnou sazbou, liší se jen barvou a titulkem.
% Rámečky jsou breakable -- delší poznámka se přelije na další stranu
% místo toho, aby přetekla přes okraj.
\newtcolorbox{lattebox}[2]{
    enhanced,
    breakable,
    colback=#1!5!white,
    colframe=#1,
    coltitle=white,
    fonttitle=\bfseries,
    title=#2,
    boxrule=0.8pt,
    arc=2pt,
    left=6pt, right=6pt, top=5pt, bottom=5pt,
}

\newenvironment{noticebox}[1][]{\begin{lattebox}{noticecolor}{#1}}{\end{lattebox}}
\newenvironment{infobox}[1][]{\begin{lattebox}{infocolor}{#1}}{\end{lattebox}}
\newenvironment{importantbox}[1][]{\begin{lattebox}{importantcolor}{#1}}{\end{lattebox}}
\newenvironment{alertbox}[1][]{\begin{lattebox}{alertcolor}{#1}}{\end{lattebox}}

% --- Rozhodovací problém ---------------------------------------------
% Ustálený zápis "vstup / otázka" používaný u rozhodovacích problémů.
% #1 = název, #2 = vstup, #3 = otázka
\newcommand{\decproblem}[3]{%
    \begin{center}
        \begin{tcolorbox}[
            enhanced, breakable, width=0.9\textwidth,
            colback=importantcolor!5!white, colframe=importantcolor,
            coltitle=white, fonttitle=\bfseries, title={#1},
        ]
            \renewcommand{\arraystretch}{1.4}%
            \begin{tabularx}{\linewidth}{@{}l X@{}}
                \textbf{Vstup:} & #2 \\
                \textbf{Otázka:} & #3 \\
            \end{tabularx}
        \end{tcolorbox}
    \end{center}%
}
